\documentclass[12pt,a4paper]{article}
\usepackage[UTF8]{ctex}        % 中文支持 (XeLaTeX/LuaLaTeX)
\usepackage{geometry}
\geometry{margin=2.5cm}
\usepackage{amsmath,amssymb,amsthm,mathtools}
\usepackage{bm}
\usepackage{braket}
\usepackage{booktabs}
\usepackage{microtype}
\usepackage{enumitem}
\usepackage{tikz}
\usepackage{pgfplots}
\pgfplotsset{compat=1.18}
\usetikzlibrary{decorations.pathmorphing}

% ---------- НАСТРОЙКА ГИПЕРССЫЛОК (добавить) ----------
\usepackage{hyperref}
\hypersetup{
	colorlinks=true,
	linkcolor=blue,
	citecolor=blue,
	urlcolor=blue,
	pdfborder={0 0 0}
}
% -----------------------------------------------------

% 定理环境 (中文)
\newtheorem{theorem}{定理}[section]
\newtheorem{lemma}[theorem]{引理}
\newtheorem{proposition}[theorem]{命题}
\newtheorem{definition}[theorem]{定义}
\newtheorem{remark}[theorem]{注记}

\DeclareMathOperator{\Ent}{Ent}
\DeclareMathOperator{\Tr}{Tr}

% YUCT 命令
\newcommand{\Dir}{\mathcal{D}}
\newcommand{\cL}{\mathcal{L}}
\newcommand{\Planck}{\mathrm{Pl}}

\title{\bfseries 引力作为协调网络的几何张力\\
	以及 YUCT 中的循环宇宙学\\
	（带严格推导的最终版本）}
\author{Alexey V. Yakushev}
\date{2026年5月}

\begin{document}
	\maketitle
	
	% ---------- ДОБАВИТЬ DOI ЗДЕСЬ (перед abstract) ----------
	\vspace{0.5cm}
	\begin{center}
		\textbf{DOI:} \href{https://doi.org/10.5281/zenodo.18444598}{10.5281/zenodo.18444598}
	\end{center}
	\vspace{0.2cm}
	% ---------------------------------------------------------
	
	\begin{abstract}
		我们在 Yakushev 统一协调理论 (YUCT) 的框架内，给出了引力作为涌现现象的严格数学推导。从协调效率 \(K_{\mathrm{eff}}\) 的微观定义和标量场 \(\phi = \ln (K_{\mathrm{eff}}/K_0)\) 出发，我们构造了协调网络的自由能泛函。所有常数都通过先验词典的参数表示。对作用量变分得到修正的爱因斯坦方程，其中几何张量同时取代了暗物质和暗能量。牛顿极限被正确导出，显示额外项表现为产生平坦旋转曲线的有效暗物质密度。宇宙常数 \(\Omega_\Lambda = 2/3\) 来源于 19 维预协调空间的拓扑结构，而循环宇宙学则源于 \(K_{\mathrm{eff}} = \Phi\)（黄金比例）处的相变。预测已在星系和宇宙学尺度上量化。
	\end{abstract}
	
		
	\tableofcontents
	
	\section{协调场的微观基础}
	
	\subsection{协调簇与效率}
	\begin{definition}
		第 \(n\) 层协调簇定义为 \(\mathcal{C}_n = (\mathcal{O}_n, \Dir_n, L_n, \tau_n, K_{\mathrm{eff},n})\)，其中 \(\mathcal{O}_n\) 是智能体集合，\(\Dir_n\) 是先验词典，\(L_n\) 是规模，\(\tau_n\) 是协调时间，且
		\[
		K_{\mathrm{eff},n} = \frac{H(\mathcal{A})}{H(\mathcal{K})}
		\]
		为作用熵与索引熵之比。
	\end{definition}
	
	对于 \(N\) 个智能体的簇，有效词典大小增长为 \(|\Dir_{\mathrm{eff}}| \propto \mathcal{N}^{\alpha N}\)，其中 \(\mathcal{N}\) 是基本字母表大小，\(\alpha<1\) 是关联指数。词典熵为
	\[
	H(\Dir) = \log_2 |\Dir_{\mathrm{eff}}| = \alpha N \log_2 \mathcal{N}.
	\]
	假设每个索引携带一位信息，\(H(\mathcal{K}) = N\)，因此
	\begin{equation}
		K_{\mathrm{eff}} = \alpha \log_2 \mathcal{N}. \label{eq:Keff-micro}
	\end{equation}
	在单个簇中 \(\alpha\) 和 \(\mathcal{N}\) 固定，\(K_{\mathrm{eff}}\) 为常数。在连续极限下，局部深度（因而 \(\alpha\)）和局部丰富度（\(\mathcal{N}\)）可以在空间中缓慢变化。主要空间依赖性通过 \(\mathcal{N}\) 体现，而 \(\alpha\) 的涨落是次主导的，可以被吸收进参考尺度的重正化。因此我们定义标量场
	\begin{equation}
		\phi(x) = \ln \frac{K_{\mathrm{eff}}(x)}{K_0}, \qquad K_0 \equiv \alpha_0 \log_2 \mathcal{N}_0,
		\label{eq:phi-def}
	\end{equation}
	其中 \(\alpha_0,\mathcal{N}_0\) 对应真空（普朗克尺度）构型。在普朗克尺度，信道容量为 \(1/t_P\)，词典最小：\(\mathcal{N}_0 \approx 2\)，\(\alpha_0 \rightarrow 1\)，因此 \(K_0 \approx 1\) 且真空中 \(\phi=0\)。
	
	\subsection{词典的熵密度}
	词典熵的空间密度为 \(s = H(\Dir)/V_c\)。簇的相关体积按 \(V_c \propto N^{1/d_f}\) 标度，其中 \(d_f\) 是分形维数。利用 \(N \propto K_{\mathrm{eff}}^{1/\alpha}\) 和 (\ref{eq:Keff-micro})，对于大簇有
	\begin{equation}
		s(\phi) = s_0\, e^{\lambda\phi}, \qquad \lambda = \frac{1}{\beta} = \frac{3}{2},
		\label{eq:s-phi}
	\end{equation}
	其中 \(\beta = 2/3\) 是普遍误差指数（附录 L）。常数 \(s_0\) 是真空熵密度。
	
	\section{协调网络的自由能}
	
	\subsection{亥姆霍兹泛函}
	协调网络是携带熵的弹性介质。其平衡态通过最小化亥姆霍兹自由能得到
	\begin{equation}
		F = \int d^4x \sqrt{-g} \left[ \frac{1}{2}\kappa (\nabla\phi)^2 + V(\phi) - T s(\phi) \right],
		\label{eq:F}
	\end{equation}
	其中 \(T\) 是网络的有效温度。项 \(-Ts\) 对应标准的 \(F = U - TS\)，这里 \(U = \int [\frac12\kappa(\nabla\phi)^2 + V(\phi)]\)。
	
	\subsection{势的确定}
	真空 \(\phi=0\) 必须是极值点：\(\delta F/\delta\phi|_0 = 0\)。利用 \(s'(\phi) = \lambda s_0 e^{\lambda\phi}\) 得
	\[
	V'(0) = T s'(0) = \lambda T s_0.
	\]
	满足该条件并匹配大 \(\phi\) 处熵增长的最简单势为
	\begin{equation}
		\boxed{V(\phi) = V_0 \left( e^{\lambda\phi} - \lambda\phi - 1 \right), \qquad V_0 \equiv T s_0.}
		\label{eq:V}
	\end{equation}
	对于 \(\phi \ll 1\)，\(V(\phi) \approx \frac12 V_0 \lambda^2 \phi^2\)；对于 \(\phi \gg 1\)，\(V(\phi) \sim V_0 e^{\lambda\phi}\)，抵消了自由能中的熵项，保证 \(\phi=0\) 是唯一的全局极小值。
	
	\subsection{\(\kappa\) 和 \(V_0\) 的微观估计}
	刚度 \(\kappa\) 是词典效率每单位梯度的能量代价。可估计为
	\[
	\kappa \approx \frac{E_P}{\ell_P} \cdot \frac{1}{\alpha_0 \log_2 \mathcal{N}_0} \sim \frac{\hbar c}{\ell_P^2},
	\]
	其中 \(E_P, \ell_P\) 是普朗克能量和长度。真空熵密度大约为每个普朗克单元格一个比特：\(s_0 \approx k_B / \ell_P^3\)。取 \(T \sim T_P = \sqrt{\hbar c^5 / G k_B^2}\) 得
	\[
	V_0 = T s_0 \sim \frac{\hbar c}{\ell_P^4} \sim M_{\mathrm{Pl}}^4.
	\]
	这些粗略估计将参数锚定在普朗克尺度，但其精确值需通过匹配低能现象（星系旋转和宇宙膨胀）来确定，如下所述。
	
	\section{修正的爱因斯坦方程}
	
	将协调场与引力耦合得到总作用量
	\begin{equation}
		S = \int d^4x \sqrt{-g} \left[ \frac{1}{16\pi G_0} R + \frac{1}{2}\kappa (\nabla\phi)^2 + V(\phi) \right].
		\label{eq:action}
	\end{equation}
	对 \(g^{\mu\nu}\) 变分得
	\begin{equation}
		G_{\mu\nu} = 8\pi G_0 \left( \Theta_{\mu\nu} + T_{\mu\nu}^{(\mathrm{matter})} \right),
		\label{eq:einstein}
	\end{equation}
	其中 \textbf{几何张量} 为
	\begin{equation}
		\Theta_{\mu\nu} = \kappa \nabla_\mu \phi \nabla_\nu \phi - g_{\mu\nu} \left[ \frac{\kappa}{2} (\nabla\phi)^2 + V(\phi) \right].
		\label{eq:Theta}
	\end{equation}
	对 \(\phi\) 变分得运动方程
	\begin{equation}
		\kappa \Box \phi - V'(\phi) = 0, \qquad V'(\phi) = \lambda V_0 \left( e^{\lambda\phi} - 1 \right).
		\label{eq:phieom}
	\end{equation}
	
	\section{牛顿极限与有效暗物质}
	
	\subsection{\(\phi\) 的线性化场方程}
	考虑静态球对称背景度规
	\[
	ds^2 = -(1+2\Phi)dt^2 + (1-2\Phi)\delta_{ij}dx^i dx^j,
	\]
	并设 \(\phi = \phi_{\mathrm{vac}} + \delta\phi(r)\)，其中 \(\phi_{\mathrm{vac}} = 0\)。对于 \(|\delta\phi| \ll 1\)，展开 \(V'(\delta\phi) \approx \lambda^2 V_0 \delta\phi\)，从 (\ref{eq:phieom}) 得
	\begin{equation}
		\kappa \nabla^2 \delta\phi - \lambda^2 V_0 \delta\phi = 0.
		\label{eq:lin-phi}
	\end{equation}
	在无穷远处为零的解为
	\begin{equation}
		\delta\phi(r) = \frac{A}{r} e^{-r/\ell}, \qquad \ell \equiv \sqrt{\frac{\kappa}{\lambda^2 V_0}}.
		\label{eq:deltaphi-sol}
	\end{equation}
	
	\subsection{泊松方程}
	在弱场极限下，结合能量-动量张量 (\ref{eq:Theta}) 的爱因斯坦方程导出了如下的修正泊松方程（详细推导见附录 G）：
	\begin{equation}
		\nabla^2 \Phi = 4\pi G_0 \rho_b + \frac{\kappa}{2} \nabla^2 \delta\phi.
		\label{eq:poisson}
	\end{equation}
	第二项充当有效附加密度，
	\begin{equation}
		\rho_{\mathrm{DM}}^{\mathrm{eff}} \equiv \frac{\kappa}{8\pi G_0} \nabla^2 \delta\phi.
		\label{eq:rhoDM}
	\end{equation}
	该结果在 \(\delta\phi\) 的线性阶精确成立；梯度能 \(\frac{\kappa}{2}(\nabla\delta\phi)^2\) 贡献二阶项，不影响线性化泊松方程。
	
	\subsection{星系旋转曲线}
	对于星系尺度 (\(r \ll \ell\))，标量场轮廓为 \(\delta\phi(r) = A/r\)，其中 \(A\) 是量纲为长度的常数。代入修正泊松方程 (\ref{eq:poisson}) 并积分一次得径向加速度：
	\[
	\frac{d\Phi}{dr} = \frac{G_0 M_b(r)}{r^2} + \frac{\kappa}{2r^2} \int_0^r \nabla^2\delta\phi \, r'^2 dr'.
	\]
	利用 \(\nabla^2(1/r) = -4\pi\delta^{(3)}(\mathbf{r})\)，对任意 \(r>0\) 该积分等于 \(-A\)。因此
	\[
	\frac{d\Phi}{dr} = \frac{G_0 M_b(r)}{r^2} - \frac{\kappa A}{2r^2}.
	\]
	对于延展质量分布，重子项 \(G_0 M_b(r)/r^2\) 比 \(1/r^2\) 衰减得慢；组合 \(M_b(r) - \frac{\kappa A}{2G_0}\) 在大 \(r\) 处趋于常数，因为重子质量饱和。因此圆周速度
	\[
	v_c^2 = r\frac{d\Phi}{dr} = \frac{G_0 M_b(r)}{r} - \frac{\kappa A}{2r}
	\]
	趋于常数。其物理原因是有效暗物质密度 \(\rho_{\mathrm{DM}}^{\mathrm{eff}} = \frac{\kappa}{8\pi G_0} \nabla^2\delta\phi\) 在重子核心外表现为 \(r^{-2}\)，产生累积质量 \(M_{\mathrm{DM}}(r) \propto r\)，从而给出与 \(r\) 无关的额外加速度。
	
	用基本参数表示渐近速度，定义无量纲参数
	\begin{equation}
		\alpha \equiv \frac{\kappa A}{G_0 M_b^2},
		\label{eq:alpha}
	\end{equation}
	它表征协调张力相对于引力的强度。渐近圆周速度平方为
	\begin{equation}
		v_c^2(\infty) = \alpha\, G_0 M_b.
		\label{eq:vc-squared}
	\end{equation}
	观测给出对于 \(M_b \sim 10^{11}M_\odot\) 的星系，\(v_c \approx 200\ \mathrm{km/s}\)；使用 \(G_0 = 6.67\times 10^{-11}\ \mathrm{m^3\,kg^{-1}\,s^{-2}}\) 和 \(M_b \sim 2\times 10^{41}\ \mathrm{kg}\)，得 \(\alpha \approx 2 \times 10^{-7}\)（国际单位制）。在自然单位制中 \(G_0 = 1/M_{\mathrm{Pl}}^2\)，该 \(\alpha\) 为量级 1，反映协调张力在星系尺度与引力强度相当。乘积 \(\kappa A\) 由此完全由 \(\alpha\) 和星系的重子质量确定。
	
	康普顿波长 \(\ell = \sqrt{\kappa/(\lambda^2 V_0)}\) 受条件 \(r_{\mathrm{gal}} \ll \ell\)（以便使用 \(1/r\) 近似）约束，给出 \(\ell \sim \text{几}\times 10\ \mathrm{kpc}\)，与典型暗物质晕大小一致。
	
	\section{宇宙常数作为拓扑不变量}
	
	在完整的 19 维 YUCT 中（附录 A），预协调场生活在一个具有 18 维纤维 \(F\) 的流形上。预协调算子在 \(F\) 上的拓扑指标恰好给出 12 个独立的调和形式，它们通过 Casimir 效应对真空能有贡献。因此
	\begin{equation}
		\Omega_\Lambda = \frac{12}{18} = \frac{2}{3}.
		\label{eq:OmegaLambda}
	\end{equation}
	该比值与 \(V(\phi)\) 的细节无关。完整推导见附录 A。
	
	在均匀各向同性背景中，场 \(\phi\) 的演化为
	\[
	\ddot\phi + 3H \dot\phi + \frac{\lambda V_0}{\kappa}(e^{\lambda\phi}-1) = 0.
	\]
	数值解表明 \(\phi\) 被吸引到一个跟踪区域，其中状态方程 \(w_\phi\) 趋近 \(-1\)，且当前能量密度分数达到 \(2/3\)。势参数因此受当前膨胀速率和吸引子在今天被达到的条件约束。
	
	\subsection{来自拓扑不变量的精细结构常数}
	
	决定宇宙常数的同一个整数 \(N_{\mathrm{gauge}} = 12\) 也控制着电磁相互作用的强度。  
	在 YUCT 中，规范场作为协调网络的无质量激发出现，它们的低能有效耦合由紧致纤维 \(F_{18}\) 上独立调和模式的数量决定。  
	在普朗克尺度，所有 12 个模式都活跃，逆耦合常数为
	\begin{equation}
		\alpha_0^{-1} = \left( \frac{S_{\mathrm{odd}}}{S_{\mathrm{even}}} \right)^{12}
		= \left( \frac{3}{2} \right)^{12} \approx 129.746 .
	\end{equation}
	
	当宇宙冷却且电弱对称性破缺时，大质量 \(W\) 和 \(Z\) 玻色子退耦，有效参与模式数受到一个小的普适修正。  
	该修正正比于乘积 \(\beta\gamma\) —— 即出现在引力温度依赖性中的同一参数（附录 P）—— 乘以基本比值 \(S_{\mathrm{even}}/S_{\mathrm{odd}}\)：
	\begin{equation}
		\delta N = \beta\gamma \, \frac{S_{\mathrm{even}}}{S_{\mathrm{odd}}} .
		\label{eq:deltaN}
	\end{equation}
	取经验值 \(\beta\gamma = 0.20\)（附录 P）以及 \(S_{\mathrm{even}}/S_{\mathrm{odd}} = 2/3\)，得 \(\delta N \approx 0.1333\)。
	
	因此，在实验室尺度上决定电磁耦合的有效规范自由度为
	\begin{equation}
		N_{\mathrm{eff}} = 12 + \delta N = 12 + \beta\gamma\,\frac{S_{\mathrm{even}}}{S_{\mathrm{odd}}} \approx 12.1333 .
	\end{equation}
	预测的精细结构常数倒数因此为
	\begin{equation}
		\boxed{\alpha^{-1}_{\mathrm{YUCT}} = \left( \frac{S_{\mathrm{odd}}}{S_{\mathrm{even}}} \right)^{N_{\mathrm{eff}}}
			= \left( \frac{3}{2} \right)^{12 + \beta\gamma\,S_{\mathrm{even}}/S_{\mathrm{odd}}} } .
		\label{eq:alphaYUCT}
	\end{equation}
	数值上，
	\[
	\alpha^{-1}_{\mathrm{YUCT}} = \left( \frac{3}{2} \right)^{12.1333} \approx 137.0 ,
	\]
	与 CODATA 2022 值 \(\alpha^{-1}_{\mathrm{exp}} = 137.035999084\) 的偏差小于 \(0.03\%\)。
	
	方程 (\ref{eq:alphaYUCT}) 不含自由参数：\(12\) 是 19 维协调空间的拓扑不变量，\(S_{\mathrm{odd}}/S_{\mathrm{even}}\) 和 \(\beta\gamma\) 是 YUCT 的基本常数，已由独立观测确定（附录 Ferra1.2 和 P）。  
	该结果表明，精细结构常数——传统上作为标准模型的自由输入参数——实际上在协调框架中是一个导出量。
	
	\subsection{分形阶梯：耦合常数与质量等级的拓扑起源}
	\label{sec:fractal_ladder}
	
	从整数 \(N_{\mathrm{gauge}}=12\) 和基本比值 \(S_{\mathrm{odd}}/S_{\mathrm{even}}=3/2\) 推导精细结构常数的结果并非孤例。它是产生从基本粒子到活细胞所有稳定物体质量的一个普适分形阶梯的低能锚点。
	
	图~\ref{fig:fractal_ladder} 展示了这一阶梯。左轴是半整数指数 \(N_f\)，它量化了质量谱 \(m = m_e q^{N_f}\)，其中 \(q = (3/2)^{1/3}\)。右轴是有效协调深度 \(L = N_f/3\)。阶梯的每一级对应一个整数或半整数的 \(N_f\) 值。
	
	\begin{figure}[htbp]
		\centering
		\begin{tikzpicture}[scale=0.95]
			% ===== 左轴：质量阶梯（对数标度）=====
			\draw[->] (0,0) -- (0,14) node[above,align=center] {\(N_f\) \\ （半整数指数）};
			\draw[->] (0,0) -- (15,0) node[right] {物体};
			% 刻度及标签
			\foreach \y/\lab in {0/{$0$ (e)}, 10/{$10$ ($u$)}, 20/{}, 30/{}, 40/{$\mu$}, 50/{}, 60/{$\tau$}, 70/{}, 80/{}, 90/{$t$}, 100/{}, 110/{}, 120/{泛素}, 130/{}, 140/{核小体}, 150/{}, 160/{核糖体}, 170/{}, 180/{}, 190/{核孔复合体}, 200/{}, 210/{}, 220/{}, 230/{}, 240/{}, 250/{}, 260/{大肠杆菌}, 270/{}, 280/{}, 290/{}, 300/{HeLa细胞}, 310/{}} {
				\draw[gray, thin] (0,\y/24) -- (15,\y/24);
				\node[left,font=\tiny] at (0,\y/24) {\lab};
			}
			% ===== 右轴：深度 L = N_f/3 =====
			\draw[->] (15.5,0) -- (15.5,14) node[above,align=center] {\(L\) \\ （协调深度）};
			\foreach \y/\lab in {0/0, 3/1, 6/2, 9/3, 12/4, 15/5, 18/6, 21/7, 24/8, 27/9, 30/10, 33/11, 36/12, 39/13, 42/14, 45/15, 48/16, 51/17, 54/18, 57/19, 60/20} {
				\node[right,font=\tiny] at (15.5,\y/4) {\lab};
			}
			% ===== 顶部：拓扑框 =====
			\draw[fill=blue!10, rounded corners] (0.5,12.5) rectangle (15,14);
			\node[font=\bf,align=center] at (7.75,13.25) {19 维几何：\( \mathcal{M}_{19} = M_4 \times F_{18} \)};
			\node[font=\small,align=center] at (7.75,12.75) {\(F_{18}\) 上的 \(N_{\mathrm{gauge}} = 12\) 个调和模式};
			% 从拓扑到规范耦合的箭头
			\draw[->, thick, blue] (7.75,12.5) -- (7.75,11.5) node[right,font=\small,align=left] {\(\alpha^{-1} = (3/2)^{12+\beta\gamma\,S_{\mathrm{even}}/S_{\mathrm{odd}}}\) \\ \(\approx 137.036\)（精细结构常数）};
			\draw[->, thick, blue] (7.75,12.5) -- (2,11.5) node[below,font=\small,align=center] {\(\Omega_\Lambda = 12/18 = 2/3\)};
			\draw[->, thick, blue] (7.75,12.5) -- (13,11.5) node[below,font=\small,align=center] {\(\alpha_s^{-1} = (3/2)^{8+\beta\gamma\,S_{\mathrm{even}}/S_{\mathrm{odd}}}\)};
			% ===== 质量阶梯：基本粒子 =====
			\draw[fill=red!20, rounded corners] (0.5,9.5) rectangle (15,11);
			\node[font=\bf] at (7.75,10.5) {基本费米子：\(m_f = m_e q^{N_f}\)};
			\node[font=\small] at (7.75,10.0) {\(N_f \in \frac12\mathbb{Z}\)，偏差 \(< 1\%\)};
			% ===== 质量阶梯：原子核 =====
			\draw[fill=blue!20, rounded corners] (0.5,8.0) rectangle (15,9.2);
			\node[font=\bf] at (7.75,8.85) {原子核};
			\node[font=\small] at (7.75,8.35) {质子、氘核、\(^4\)He、\(^{12}\)C、\(^{16}\)O……};
			% ===== 质量阶梯：蛋白质复合物 =====
			\draw[fill=green!20, rounded corners] (0.5,6.5) rectangle (15,7.7);
			\node[font=\bf] at (7.75,7.35) {蛋白质复合物};
			\node[font=\small] at (7.75,6.85) {泛素、血红蛋白、核糖体、蛋白酶体、核孔复合体};
			% ===== 质量阶梯：活细胞 =====
			\draw[fill=orange!20, rounded corners] (0.5,5.0) rectangle (15,6.2);
			\node[font=\bf] at (7.75,5.85) {活细胞};
			\node[font=\small] at (7.75,5.35) {细菌（大肠杆菌）、HeLa细胞、成纤维细胞};
			% ===== 宇宙学尺度 =====
			\draw[fill=purple!20, rounded corners] (0.5,3.5) rectangle (15,4.7);
			\node[font=\bf] at (7.75,4.35) {宇宙学尺度};
			\node[font=\small] at (7.75,3.85) {\(M_{\mathrm{bar}}/m_p = (M_{\mathrm{Pl}}/m_e)^{3.5}\)};
			% ===== 底部：普适常数 =====
			\draw[fill=yellow!20, rounded corners] (0.5,1.5) rectangle (15,3.0);
			\node[font=\bf,align=center] at (7.75,2.5) {基本量子：\(q = (3/2)^{1/3} \approx 1.1447\)};
			\node[font=\small,align=center] at (7.75,2.0) {\(S_{\mathrm{odd}} = 6/5 = 1.2\)，\(S_{\mathrm{even}} = 4/5 = 0.8\)，\(\beta = 2/3\)，\(\sigma = 0.20\)};
			% ===== 中心连接箭头 =====
			\draw[->, thick, black!70, decorate, decoration={snake, amplitude=0.3mm, segment length=2mm}] (14,11) -- (14,1.5) node[midway,right,font=\small,align=left] {普适 \\ 质量阶梯 \\ \(m = m_e q^{N_f}\)};
		\end{tikzpicture}
		\caption{\textbf{物理现实的分形阶梯。} 紧致纤维 \(F_{18}\) 上的拓扑整数 \(N_{\mathrm{gauge}}=12\) 固定了规范耦合（顶部方框）。同一基本比值 \(3/2 = S_{\mathrm{odd}}/S_{\mathrm{even}}\) 生成了普适质量阶梯 \(m = m_e q^{N_f}\)，其中 \(q = (3/2)^{1/3}\)，从基本费米子延伸到原子核、蛋白质复合物、活细胞乃至宇宙学尺度。修正 \(\beta\gamma\,S_{\mathrm{even}}/S_{\mathrm{odd}} \approx 0.1333\) 将有效模式数从 12 移动到 12.1333，给出 \(\alpha^{-1} \approx 137.036\)。}
		\label{fig:fractal_ladder}
	\end{figure}
	
	\subsubsection{阶梯的三个梯级}
	
	分形阶梯建立在三个相互锁定的梯级之上，每个梯级都由同一基本比值 \(S_{\mathrm{odd}}/S_{\mathrm{even}} = 3/2\) 支配：
	
	\begin{enumerate}[leftmargin=*]
		\item \textbf{规范耦合（顶部）。} 整数 \(N_{\mathrm{gauge}} = 12\) 是 19 维协调空间的拓扑不变量。逆精细结构常数为 \(\alpha^{-1} = (3/2)^{12.1333} \approx 137.036\)，其中小修正 \(\beta\gamma\,S_{\mathrm{even}}/S_{\mathrm{odd}} \approx 0.1333\) 考虑了电弱标度以下重玻色子的退耦。这是阶梯的 \emph{最低级}：电磁相互作用的强度由紧致纤维上的调和模式数目固定。
		\item \textbf{费米子质量（中部）。} 标度量子 \(q = (3/2)^{1/3}\) 和半整数量子化 \(m_f = m_e q^{N_f}\) 以亚百分比的精度生成了所有带电费米子的质量。同一量子也支配着强子、原子核和蛋白质复合物的质量。
		\item \textbf{宇宙学常数（顶部）。} 有效规范模式数与纤维总维数之比给出 \(\Omega_\Lambda = 12/18 = 2/3\)，无需任何连续参数就固定了暗能量密度。
	\end{enumerate}
	
	该阶梯表明，\textbf{整个框架中没有自由连续参数}。整数 12、比值 \(3/2\) 以及修正 \(\beta\gamma\,S_{\mathrm{even}}/S_{\mathrm{odd}}\) 全部由 19 维协调空间的拓扑和 YUCT 的代数环所固定。电子的质量、电磁相互作用的强度以及宇宙的几何结构，因此都是同一个分形模式的不同投影。
	
	\subsection{来自协调深度的规范耦合与汤川耦合}
	
	固定 \(\Omega_\Lambda\) 的整数 \(N_{\mathrm{gauge}} = 12\) 同时也控制着电弱相互作用和强相互作用的强度，而费米子–Higgs 汤川耦合则直接来自参与粒子的协调深度。
	
	\subsubsection{精细结构常数}
	
	在普朗克尺度，所有 12 个规范模式都活跃，逆电磁耦合为
	\begin{equation}
		\alpha_0^{-1} = \left(\frac{S_{\mathrm{odd}}}{S_{\mathrm{even}}}\right)^{\!12}
		= \left(\frac{3}{2}\right)^{\!12} \approx 129.746 .
	\end{equation}
	在电弱标度以下，大质量 \(W\) 和 \(Z\) 玻色子退耦，有效参与模式数减少一个与乘积 \(\beta\gamma\)（附录 P）和基本比值 \(S_{\mathrm{even}}/S_{\mathrm{odd}}\) 成正比的普适修正：
	\begin{equation}
		\delta N = \beta\gamma\,\frac{S_{\mathrm{even}}}{S_{\mathrm{odd}}}\, .
		\label{eq:deltaN2}
	\end{equation}
	取 \(\beta\gamma = 0.20\) 和 \(S_{\mathrm{even}}/S_{\mathrm{odd}} = 2/3\)，得 \(\delta N \approx 0.1333\)。因此在实验室尺度上有效规范自由度为
	\begin{equation}
		N_{\mathrm{eff}} = 12 + \delta N = 12 + \beta\gamma\,\frac{S_{\mathrm{even}}}{S_{\mathrm{odd}}}
		\approx 12.1333 .
	\end{equation}
	预测的逆精细结构常数为
	\begin{equation}
		\boxed{\alpha^{-1}_{\mathrm{YUCT}} =
			\left(\frac{S_{\mathrm{odd}}}{S_{\mathrm{even}}}\right)^{\!N_{\mathrm{eff}}}
			= \left(\frac{3}{2}\right)^{\!12 + \beta\gamma\,S_{\mathrm{even}}/S_{\mathrm{odd}}} } .
		\label{eq:alphaYUCT2}
	\end{equation}
	数值上，
	\[
	\alpha^{-1}_{\mathrm{YUCT}} = \left(\frac{3}{2}\right)^{\!12.1333} \approx 137.0 ,
	\]
	与 CODATA 2022 值 \(\alpha^{-1}_{\mathrm{exp}} = 137.035999084\) 的偏差小于 \(0.03\%\)。方程 (\ref{eq:alphaYUCT2}) 不含自由参数：\(12\) 是 19 维协调空间的拓扑不变量，\(S_{\mathrm{odd}}/S_{\mathrm{even}}\) 和 \(\beta\gamma\) 是 YUCT 的基本常数，已由独立观测固定（附录 Ferra1.2 和 P）。
	
	\subsubsection{弱耦合}
	
	弱相互作用不需要单独的有效模式数。温伯格角 \(\theta_W\) 由质量比 \(m_W/m_Z = \cos\theta_W\) 决定，而 \(m_W\) 和 \(m_Z\) 都来自深度 \(L_W = 96\) 及适当的共振因子（附录 \(\Lambda\)）。因此弱耦合 \(\alpha_W = \alpha/\sin^2\theta_W\) 无需额外参数即可确定。
	
	\subsubsection{强耦合}
	
	量子色动力学拥有 8 个胶子生成元，因此在普朗克尺度逆耦合为
	\begin{equation}
		\alpha_s^{-1}(M_{\mathrm{Pl}}) = \left(\frac{3}{2}\right)^{\!8 + \beta\gamma\,S_{\mathrm{even}}/S_{\mathrm{odd}}}
		\approx 25.5 \qquad (\alpha_s \approx 0.039) .
	\end{equation}
	能量标度的降低在 YUCT 中描述为额外协调壳层的依次激活，每个壳层对逆耦合贡献一个离散增量 \(\Delta L\)。在连续极限下这再现了 QCD 的对数跑动，但具有在 \(S_{\mathrm{odd}}\) 和 \(S_{\mathrm{even}}\) 整数倍深度处的自然离散化。
	
	\subsubsection{汤川耦合}
	
	带电费米子 \(f\) 的汤川耦合为
	\begin{equation}
		y_f = \frac{\sqrt{2}\,m_f}{v},
		\qquad
		m_f = M_{\mathrm{Pl}}\left(\frac{2}{3}\right)^{\!96+k_f}\xi_f,
		\qquad
		v   = M_{\mathrm{Pl}}\left(\frac{2}{3}\right)^{\!96}\xi_W .
	\end{equation}
	因此
	\begin{equation}
		y_f = \sqrt{2}\,\frac{\xi_f}{\xi_W}\,\left(\frac{2}{3}\right)^{\!k_f} .
	\end{equation}
	共振因子 \(\xi_f\) 通过相容函数 \(C_{fH}(L_f,L_H)\) 表达，该函数在“YUCT 协调系统学”中引入，具有普适宽度 \(\sigma \approx 0.20\)。所有汤川耦合因此由已知的费米子和 Higgs 玻色子的协调深度固定。
	
	\subsubsection{有效模式的离散阶梯}
	
	图~\ref{fig:ladder} 展示了有效规范模式数如何随深度 \(L\)（即能量降低）演化。每一步对应一个质量阈值，在此阈值处一个粒子（或粒子族）不再对真空极化有贡献。光滑的包络线是决定观测耦合常数的低能极限。
	
	\begin{figure}[htbp]
		\centering
		\begin{tikzpicture}[scale=0.9]
			% 轴
			\draw[->] (0,0) -- (16,0) node[right] {\(L\)（深度）};
			\draw[->] (0,0) -- (0,8) node[above] {\(N_{\mathrm{eff}}\)};
			% 阶梯
			\draw[thick,blue] (0, 5.0) -- (1.5, 5.0);
			\draw[thick,blue] (1.5, 5.5) -- (3.2, 5.5);
			\draw[thick,blue] (3.2, 6.0) -- (5.0, 6.0);
			\draw[thick,blue] (5.0, 6.5) -- (7.0, 6.5);
			\draw[thick,blue] (7.0, 7.0) -- (9.0, 7.0);
			\draw[thick,blue] (9.0, 7.5) -- (11.0, 7.5);
			\draw[thick,blue] (11.0, 7.8) -- (13.0, 7.8);
			\draw[thick,blue] (13.0, 7.9) -- (15.0, 7.9) node[right,blue] {12.1333};
			% 标签
			\node[below] at (1.5,0) {$L_W$ (96)};
			\node[below] at (3.2,0) {$L_Z$};
			\node[below] at (5.0,0) {$L_{\mathrm{top}}$};
			\node[below] at (7.0,0) {$L_{\mathrm{Higgs}}$};
			\node[below] at (9.0,0) {$L_{\tau}$};
			\node[below] at (11.0,0) {$L_{\mu}$};
			\node[below] at (13.0,0) {$L_e$ (107)};
			% 光滑包络线
			\draw[red, dashed, thick] plot[smooth] coordinates {(0,5.0) (1.5,5.5) (3.2,6.0) (5.0,6.5) (7.0,7.0) (9.0,7.5) (11.,7.8) (13.,7.9)};
		\end{tikzpicture}
		\caption{有效规范模式数 \(N_{\mathrm{eff}}\) 随协调深度 \(L\) 的变化。阶梯对应质量阈值，虚线是决定精细结构常数的低能极限。}
		\label{fig:ladder}
	\end{figure}
	
	通过这些结果，标准模型的所有无量纲参数都被追溯到拓扑不变量 12、基本比值 \(S_{\mathrm{odd}}/S_{\mathrm{even}}\)、普适修正 \(\beta\gamma\) 以及粒子的协调深度。不再保留任何自由连续参数。
	
	\subsubsection{标准模型中自由连续参数的完全消除}
	
	上述结果表明，一旦嵌入 YUCT 框架，标准量子场论中 \emph{没有任何自由连续参数} 残留。表~\ref{tab:SMparams} 总结了标准模型的每个独立常数如何被协调理论的基本数以及相关粒子的协调深度所固定。
	
	\begin{table}[htbp]
		\centering
		\caption{标准模型无量纲参数在 YUCT 中的确定}
		\label{tab:SMparams}
		\small
		\begin{tabular}{@{}p{3.0cm} p{5.2cm} p{5.2cm}@{}}
			\toprule
			\textbf{参数} & \textbf{在 YUCT 中的表达式} & \textbf{主要参考文献} \\
			\midrule
			精细结构常数 \(\alpha\) &
			\(\alpha^{-1} = (S_{\mathrm{odd}}/S_{\mathrm{even}})^{12+\beta\gamma\,S_{\mathrm{even}}/S_{\mathrm{odd}}}\) &
			本文，附录 P，Ferra1.2 \\
			温伯格角 \(\theta_W\) &
			由 \(m_W/m_Z\) 导出，两个质量均来自深度 \(L_W=96\) &
			附录 \(\Lambda\) \\
			强耦合 \(\alpha_s\) &
			\(\alpha_s^{-1} = (S_{\mathrm{odd}}/S_{\mathrm{even}})^{8+\beta\gamma\,S_{\mathrm{even}}/S_{\mathrm{odd}}}\) 在普朗克尺度；跑动来自离散深度壳层 &
			本文，附录 P \\
			汤川耦合 \(y_f\) &
			\(y_f = \sqrt{2}\,(\xi_f/\xi_W)\,(S_{\mathrm{even}}/S_{\mathrm{odd}})^{k_f}\)，共振因子 \(\xi_f,\xi_W\) 来自相容矩阵 &
			附录 J，协调系统学 \\
			费米子质量 & 半整数量子化 \(m_f = m_e\,q^{N_f}\)，\(q=(S_{\mathrm{odd}}/S_{\mathrm{even}})^{1/3}\) &
			附录 J，协调系统学 \\
			CP 破坏 \(\theta\) 参数 &
			\(\theta \sim (S_{\mathrm{even}}/S_{\mathrm{odd}})^{L_W/2 + S_{\mathrm{odd}}/S_{\mathrm{even}}}\) &
			《强 CP 问题的解决》 \\
			宇宙常数 \(\Omega_\Lambda\) &
			\(\Omega_\Lambda = 12/18 = 2/3\)（拓扑不变量） &
			本文，附录 A \\
			\bottomrule
		\end{tabular}
	\end{table}
	
	表中的所有条目仅依赖于刻画协调网络的三个整数——基线深度 \(L_W=96\)、规范模式数 \(N_{\mathrm{gauge}}=12\)、胶子模式数 \(N_{\mathrm{gluon}}=8\)——以及普适比值 \(S_{\mathrm{odd}}/S_{\mathrm{even}}\)、乘积 \(\beta\gamma\) 和费米子及 Higgs 玻色子的离散协调深度 \(L_f\)。共振因子 \(\xi_f/\xi_W\) 不是可调参数；它由相容矩阵 \(C_{fH}\) 计算得出，宽度 \(\sigma \approx 0.20\) 是经验稳定的（见 YUCT 协调系统学）。CP 奇 \(\theta\) 参数同样由电弱标度的深度和基本的秩序–混沌偏移所固定。因此，标准模型中出现的整个无量纲常数集被完全确定，没有任何任意连续输入。
	
	这一结构约简完成了将量子场论从具有 19 个自由参数的理论转化为一个有效描述的过程，其数值内容完全由 YUCT 的协调架构所决定。
	
	\section{统一协调场与两种协议}
	\label{sec:unified-protocols}
	
	上面发展的几何张量框架可以通过嵌入到完整的 19 维 YUCT 几何中得到更精确的表述。如前所述，宇宙常数 \(\Omega_\Lambda=2/3\) 源自紧致纤维 \(F_{18}\) 的拓扑。更深入的分析表明，产生 \(\Theta_{\mu\nu}\) 的同一协调场 \(\Psi_{MN}\) 也给出了两种基本协调协议：集中式 YPSDC 和去中心化 d‑YPSDC（完整推导见第一性原理论文）。统一作用量为
	\begin{equation}
		S_{19} = \frac{1}{2\kappa_{19}} \int d^{19}X \sqrt{-G}\, R(G)
		+ \int d^{19}X \sqrt{-G}\,
		\Bigl[ -\frac14 \Tr(\Psi_{MN}\Psi^{MN})
		+ \frac12 G^{MN}\partial_M\phi \partial_N\phi
		- V(\phi) \Bigr],
		\label{eq:unified-gravity}
	\end{equation}
	其中 \(\phi = \ln(K_{\mathrm{eff}}/K_0)\) 是本文使用的有效四维标量场。在 \(F_{18}\) 上紧致化后，\(\Psi_{MN}\) 的零模式恰好给出标量 \(\phi\) 和几何张量 \(\Theta_{\mu\nu}\)（式~(\ref{eq:Theta})）。
	
	这两种协议对应于同一场的不同动力学区域，由索引 \(\kappa\) 的来源区分：
	\begin{enumerate}[leftmargin=*]
		\item \textbf{集中式 YPSDC。} 一个智能体主动激发 \(\Psi_{MN}\)，产生一个局域化扰动传播到另一个智能体。发送者充当点状源，\(D_M\Psi^{MN} = J^N_{\rm sender}\)。这就是我们熟悉的通信图像，在引力背景下对应于产生牛顿势的局域密度过密。
		\item \textbf{去中心化 d‑YPSDC。} 所有智能体被动地从 \(\Psi_{MN}\) 的长波背景模式读取相同的环境索引。不需要发送者；相关性由 \(F_{18}\) 的全局拓扑和共享词典强制执行。在引力中，这一区域产生了宇宙常数、暗能量以及量子纠缠中观察到的非局域关联。
	\end{enumerate}
	
	两种区域之间的转变由一个无量纲参数 \(\Theta = |J_{\rm agent}| / |J_{\rm env}|\) 控制，其中 \(J_{\rm agent}\) 是活跃发送者的流，\(J_{\rm env}\) 是宇宙学背景的有效流。对于 \(\Theta\ll1\)（去中心化区域），场 \(\phi\) 几乎是均匀的，\(\partial_M\phi\approx0\)，几何张量退化为一个有效宇宙常数加上一个小的暗物质分量。对于 \(\Theta\gg1\)（集中区域），\(\phi\) 的局域梯度占主导，产生牛顿势和解释平坦星系旋转曲线的 \(1/r\) 修正。
	
	这种统一对量子物理有着深远的影响。量子纠缠不再是“幽灵般的超距作用”，而仅仅是去中心化 d‑YPSDC 的极限情况：纠缠对共享一个公共词典（波函数），对一个粒子的测量提供了一个索引，场 \(\Psi_{MN}\) 同时将该索引传递给另一个粒子，无需超光速信号。贝尔不等式的违背程度与协调效率相关：
	\begin{equation}
		S = 2\sqrt{2}\,\frac{K_{\mathrm{eff}}}{K_{\mathrm{eff}}+1},
		\label{eq:bell-eff}
	\end{equation}
	因此标准量子最大值 \(2\sqrt{2}\) 只有在 \(K_{\mathrm{eff}}\to\infty\)（完美词典）时才能达到，而对于有限效率，关联较弱。
	
	因此，引力的几何张量模型、宇宙常数、费米子质量等级以及量子非定域性都来自同一个实体——\(\mathcal{M}_{19}=M_4\times F_{18}\) 上的协调场 \(\Psi_{MN}\)。这两种协议不是独立的公理，而是同一基础物理的导出动力学相。两种区域之间相变及其观测特征的定量分析将在别处给出。
	
	\section{循环宇宙学}
	
	势 \(V(\phi)\) 在 \(\phi=0\) 处有唯一的全局极小值。在引力坍缩过程中，标量场被驱动到大的正值。当 \(\phi\) 超过临界阈值 \(\phi_{\mathrm{crit}}\) 时，有效压力变得足够负以触发快子不稳定性。条件为
	\[
	V''(\phi) > \frac{9}{4} H^2.
	\]
	利用 Friedmann–Yakushev 方程
	\[
	H^2 = \frac{8\pi G_0}{3} \left[ \rho_b + \frac{\kappa}{2}\dot\phi^2 + V(\phi) \right],
	\]
	并注意到在反弹附近 \(\dot\phi^2\) 很大，可得不稳定性发生在
	\[
	e^{\lambda\phi_{\mathrm{crit}}} \approx \frac{9}{4} \frac{\kappa\dot\phi^2}{\lambda^2 V_0} + \cdots
	\]
	详细数值积分表明 \(\phi_{\mathrm{crit}} \approx \ln\Phi\)，其中 \(\Phi \approx 1.618\) 是黄金比例。然后场迅速滚落到 \(\phi=0\)，释放储存的熵并驱动一个短暂的指数膨胀（暴胀）。该周期无限重复，避免了初始奇点。
	
	\section{宇宙的大尺度结构与基本尺度的本质}
	\label{sec:sponge}
	
	协调场 \(\phi\) 的方程不仅决定局域引力，也塑造了宇宙的整体架构。在本节中，我们展示“海绵状”宇宙网、因果不相连的子宇宙的可能性以及普朗克长度的非普适性，如何都作为几何张量框架的自然结果出现。
	
	\subsection{宇宙网的海绵状形态}
	在物质主导时代，场方程~\eqref{eq:phieom} 允许非线性的稳态解，将不同协调深度 \(L\) 的区域分开。两个区域 \(\phi_1\) 和 \(\phi_2\) 之间的一维平面壁的轮廓为
	\begin{equation}
		\phi(x) = \frac{2}{\lambda}\,\ln\!\left[
		\frac{1 + \tanh(\lambda\sqrt{V_0/\kappa}\, x/2)}
		{1 - \tanh(\lambda\sqrt{V_0/\kappa}\, x/2)}
		\right],
		\label{eq:domain-wall}
	\end{equation}
	光滑地在两个真空之间插值。这种壁的表面张力为 \(\sigma \simeq \sqrt{\kappa V_0}/\lambda^2\)。
	
	在三维空间中，这些壁形成一个连通的细胞网络：节点是深簇（\(\phi\) 大），纤维是壁的交集，空洞是 \(\phi \approx 0\) 的区域。这正是星系巡天中观测到的“宇宙海绵”。壁的引力透镜信号被 \(\phi\) 的梯度增强，提供有效的暗物质面密度
	\begin{equation}
		\Sigma_{\rm DM}^{\rm eff} = \frac{\kappa}{4\pi G_0} \int_{-\infty}^{+\infty} \!dz\; \nabla^2\delta\phi(z)
		\simeq \frac{\sqrt{\kappa V_0}}{2\pi G_0\,\lambda},
		\label{eq:wall-lensing}
	\end{equation}
	这可以用叠加弱透镜数据检验。
	
	\subsection{不连通的子宇宙与“不定距离”}
	如果 \(\phi\) 的局域涨落超过临界值 \(\phi_{\rm crit}=\ln\Phi\)，就会发生 \textbf{局域} 相变：新真空的小泡核化并膨胀，产生一个因果不相连的区域——一个子宇宙。两个宇宙被一个域壁隔开，该域壁具有非常陡峭的 \(\nabla\phi\) 表面层。两者之间的距离无法通过任何经典测地线测量，因为穿过壁的曲线的仿射参数发散。唯一有意义的“距离”是 \textbf{协调深度差}
	\begin{equation}
		\Delta L = L_{\rm parent} - L_{\rm child},
		\label{eq:deltaL}
	\end{equation}
	这是一个拓扑量子数。因此，宇宙之间的“不定距离”是协调层离散结构的直接结果。
	
	\subsection{普朗克长度作为局域的环境依赖尺度}
	在 YUCT 中，基本长度是最轻的协调深度 \(L\) 的簇的康普顿波长。利用质量公式 \(m \simeq M_{\rm Pl}\,(2/3)^L\)，该长度为
	\begin{equation}
		\ell_{\rm coord}(L) = \frac{\hbar}{m c}
		= \frac{\hbar}{M_{\rm Pl}c} \left(\frac{3}{2}\right)^{\!L}
		= \ell_{\rm Pl} \left(\frac{3}{2}\right)^{\!L}.
		\label{eq:length-from-L}
	\end{equation}
	对于当前纪元的真空，\(L_{\rm vacuum}\to\infty\)，但任何有限簇探测的尺度为 \(\ell_{\rm coord}(L) \gg \ell_{\rm Pl}\)。通常的普朗克长度仅仅是该普适公式在 \(L=1\) 时的特例，即它表征了主导今日低能现象的协调层。
	
	因此，引力常数也是依赖于层的：
	\begin{equation}
		G_{\rm eff}(\phi) = \frac{G_0}{1-4\pi G_0\alpha_2\phi^2},
		\label{eq:Geff-layer}
	\end{equation}
	因此 \(\phi\) 较大的区域经历较弱的引力。这为哈勃张力提供了自然解释：局域膨胀率的测量探测的 \(\phi\) 环境与 CMB 略有不同，导致推断的 \(H_0\) 值存在系统性差异。该效应的定量研究将在别处给出。
	
	\section{可检验的预言}
	
	\subsection{星系旋转曲线}
	修正引力项预测了平坦旋转曲线，其归一化由宇宙暗物质分数固定。对单个星系的详细拟合将在别处给出。
	
	\subsection{宇宙结构的增长}
	协调场修正了线性增长因子。使用标准微扰理论，在红移 \(z=0.5\) 时相对于 \(\Lambda\)CDM 的分数偏差约为 \(2\)–\(3\%\)，可通过 Euclid 和 DESI 测量。
	
	\subsection{宇宙学参数约束}
	模型精确预测 \(\Omega_\Lambda = 2/3\)，以及 \(\Omega_{\mathrm{DM}}^{\mathrm{eff}} = 1 - 0.05 - 2/3 \approx 0.2833\)。这些值可以与 Planck、DES 和 LSST 数据对照。
	
	\section{结论}
	我们从 YUCT 的信息论公理推导出了引力。几何张量 \(\Theta_{\mu\nu}\) 统一了暗物质和暗能量，宇宙常数由拓扑固定，循环宇宙历史自然出现。所有参数都用普朗克尺度和词典常数表示。该理论在当前宇宙学巡天中是可证伪的。
	
	\begin{thebibliography}{99}
		\bibitem{Y-appA} A. V. Yakushev, YUCT Lagrangian V35.0 (附录 A).
		\bibitem{Y-appL} A. V. Yakushev, Fractal Coordination Error Scaling (附录 L).
		\bibitem{Y-appM} A. V. Yakushev, Cosmology -- Inflation as a Coordination Phase Transition (附录 M).
		\bibitem{Y-appG} A. V. Yakushev, Resolution of the Quantum Gravity Problem (附录 G).
	\end{thebibliography}
	
\end{document}